\chapter{Dynamic Benchmarking}
\label{ch:benchmark}

\section{Why a standard test is needed}

Chapter~\ref{ch:mppt} ended with four trackers separated by less than
\SI{0.5}{\percent} on a step test, in an order that turns out not to survive a
more demanding test. This is the normal situation in MPPT work: results depend
heavily on the test signal, so comparisons are only meaningful against an agreed
one.

The European standard EN 50530~\cite{en50530} defines such a test. It specifies
trapezoidal irradiance profiles, ramping linearly between two levels, dwelling,
ramping back, and repeating, with prescribed ramp rates. Two classes are
defined, a low-to-medium class between \SI{10}{\percent} and \SI{50}{\percent}
of nominal irradiance with gentle slopes, and a medium-to-high class between
\SI{30}{\percent} and \SI{100}{\percent} with steeper slopes. The figure of
merit is the dynamic MPPT efficiency of~\eqref{eq:eta}.

\section{Test profiles used}

The profiles used here follow the medium-to-high class: trapezoidal ramps
between \SI{300}{} and \SI{1000}{\watt\per\square\metre} at \SI{25}{\celsius},
with a \SI{1}{\second} dwell at each plateau and a \SI{1}{\second} settling
period before the first ramp. Four ramp rates were tested: \SI{25}{},
\SI{50}{}, \SI{100}{} and \SI{200}{\watt\per\square\metre\per\second}.

Two deviations from the standard are stated explicitly for honesty. First, one
up-down cycle is used per slope rather than the full repetition count, which
keeps each simulation to a practical length while still exercising both ramp
directions. Second, the fastest rate of \SI{200}{\watt\per\square\metre\per\second}
is above the standard's range, and was included deliberately to push the
trackers past the point where they can keep up. Neither deviation affects the
comparison between trackers, since all four see identical profiles.

\section{Computing the efficiency honestly}

The denominator of~\eqref{eq:eta} requires the true available power at every
instant. Evaluating the exact maximum power point at each of the roughly
\num{500000} simulation samples would be prohibitively slow, so a lookup is
built once: the exact MPP is solved on a grid of \num{240} irradiance values
spanning the profile range, and the result is interpolated with a shape
preserving cubic. Because $\Pmp(G)$ is smooth and nearly linear over this range,
the interpolation error is far below the differences being measured.

The score window excludes the first second, so that the initial acquisition
transient, which is an artefact of the arbitrary starting point, does not
contaminate the tracking measurement.

\section{Results}

\begin{table}[htbp]
\centering
\caption{Dynamic MPPT efficiency (\si{\percent}) on trapezoidal ramps between \SI{300}{} and \SI{1000}{\watt\per\square\metre}.}
\label{tab:dyneta}
\begin{tabular}{@{}lrrrr@{}}
\toprule
\textbf{Tracker} & \SI{25}{} & \SI{50}{} & \SI{100}{} & \SI{200}{\watt\per\square\metre\per\second} \\
\midrule
P\&O, fixed step        & 99.952 & 99.930 & 99.871 & 99.483 \\
P\&O, variable step     & 99.836 & 98.726 & \textbf{93.141} & 99.652 \\
Incremental conductance & 99.950 & 99.926 & 99.871 & 99.458 \\
Extremum seeking        & 99.851 & 99.851 & 99.851 & \textbf{99.850} \\
\bottomrule
\end{tabular}
\end{table}

\begin{table}[htbp]
\centering
\caption{Worst instantaneous power shortfall (\si{\watt}) on the same tests, against a peak available power of \SI{2402}{\watt}.}
\label{tab:dyndef}
\begin{tabular}{@{}lrrrr@{}}
\toprule
\textbf{Tracker} & \SI{25}{} & \SI{50}{} & \SI{100}{} & \SI{200}{\watt\per\square\metre\per\second} \\
\midrule
P\&O, fixed step        & 6.24  & 8.21   & 23.37  & 116.17 \\
P\&O, variable step     & 43.58 & 240.29 & 704.85 & 40.50 \\
Incremental conductance & 6.84  & 9.15   & 17.93  & 129.26 \\
Extremum seeking        & \textbf{6.43} & \textbf{6.63} & \textbf{7.04} & \textbf{7.98} \\
\bottomrule
\end{tabular}
\end{table}

\begin{figure}[htbp]
\centering
\bigfig{s05_dynamic_benchmark.png}
\caption[Dynamic MPPT benchmark]{Dynamic benchmark results. Top left, efficiency against ramp rate; extremum seeking is flat while the hill-climbers degrade, and the variable-step P\&O collapses at \SI{100}{\watt\per\square\metre\per\second}. Top right, the test profile. Bottom left, voltage tracking, where the runaway excursion of the variable-step tracker down past \SI{90}{\volt} is plainly visible. Bottom right, instantaneous power shortfall.}
\label{fig:benchmark}
\end{figure}

\section{Discussion}

\subsection{Extremum seeking is essentially rate-independent}

The most striking result in Table~\ref{tab:dyneta} is the ESC row: \SI{99.851}{},
\SI{99.851}{}, \SI{99.851}{}, \SI{99.850}{\percent}. An eightfold change in ramp
rate produces a change of \SI{0.001}{\percent}. Its worst instantaneous
shortfall rises only from \SI{6.4}{} to \SI{8.0}{\watt}, whereas the
hill-climbers rise to \SI{116}{} and \SI{129}{\watt}.

The reason is structural. ESC estimates the gradient continuously from a dither
at a known frequency, and the washout filter explicitly rejects the slowly
varying component of power. A slow irradiance ramp is precisely such a slowly
varying component, so ESC is largely blind to it by construction. The
hill-climbers, by contrast, sample power at \SI{50}{\hertz} and cannot
distinguish a power change caused by their own perturbation from one caused by
the ramp.

This is a genuinely control-theoretic advantage: frequency separation between
the probe and the disturbance, rather than a better search heuristic.

\subsection{The cost of that immunity}

ESC pays for this with its dither. Under constant irradiance it scores
\SI{99.55}{\percent} against \SI{99.94}{\percent} for fixed-step P\&O
(Table~\ref{tab:stepcompare}), because the \SI{6}{\volt} peak-to-peak excursion
sits permanently on a curved P--V characteristic and the average power on an arc
is below the power at its centre. The optimal choice therefore depends on the
site: ESC for locations with frequent fast irradiance transients, P\&O for
stable clear-sky conditions.

\subsection{Fixed-step P\&O and INC remain nearly identical}

The two classical hill-climbers agree to within \SI{0.03}{\percent} at every
ramp rate. As argued in Chapter~\ref{ch:mppt} this is expected: with the same
step and sample period they compute the sign of the same quantity by two routes.
INC's advantages are real but appear only in the presence of measurement noise
and quantisation, which this idealised simulation does not include. Adding
sensor noise models is identified as future work in
Chapter~\ref{ch:conclusions}.

\subsection{The variable-step failure}

The variable-step P\&O, which was the \emph{best} tracker under the step test,
is by far the worst here, dropping to \SI{93.14}{\percent} at
\SI{100}{\watt\per\square\metre\per\second} with a peak shortfall of
\SI{705}{\watt}, nearly a third of the available power.

The mechanism is the drift problem of Section~\ref{sec:drift} amplified by the
adaptive rule. The step size is made proportional to $|\Delta P / \Delta V|$,
which is intended to shrink the step near the flat peak and enlarge it on the
steep flanks. During a rising ramp, though, the algorithm locks into the wrong
direction and starts moving away from the peak. As it does so
$|\mathrm{d}P/\mathrm{d}V|$ \emph{grows}, so the adaptive rule enlarges the
step, which moves it further away faster. The feedback is positive. The bottom
left panel of Figure~\ref{fig:benchmark} shows the excursion running from
\SI{158}{\volt} down past \SI{90}{\volt} before the ramp ends and the algorithm
recovers.

The non-monotonic efficiency across ramp rates, worst at \SI{100}{} but
acceptable at \SI{200}{\watt\per\square\metre\per\second}, is consistent with
this explanation: whether a runaway happens depends on the phase relationship
between the perturbation sequence and the ramp, not smoothly on the ramp rate.

This result is reported rather than tuned away because it is instructive. It
demonstrates that an algorithm modification which improves the usual benchmark
can introduce a serious failure mode elsewhere, and that positive feedback
inside an adaptation rule deserves the same stability scrutiny as the main
control loop. The standard remedy in the literature is to add an irradiance
change detector, which tests whether a large $\Delta P$ was accompanied by a
correspondingly small $\Delta V$ and, if so, suppresses the direction reversal.
Implementing and validating such a drift-free variant is the most immediately
valuable extension to this project.
